% ===================== CHAPTER 6 — CONCLUSION AND FUTURE WORK =====================
\section{Conclusion}
% GUIDE: compare results with similar products; state what was/was not achieved;
% identify the main contributions; note lessons learned.
This thesis has designed and implemented a web-based supply-chain management
system for small and medium-sized distributors. Compared with broad ERP products
such as Odoo and SAP Business One, the system is narrower in scope but more
focused on the workflow selected for the thesis: purchase approval, goods
receipt, lot-based inventory, sales order approval, stock reservation, FEFO
goods issue, delivery execution, invoice generation, dashboard reporting, and
role-based access control. Compared with retail-oriented systems such as
KiotViet, Sapo, and MISA, the prototype emphasizes distributor workflow
integration rather than point-of-sale operation.

The main achieved result is a working application structure that connects the
procure-to-stock and order-to-cash processes. In terms of scale, the
requirements were modelled as nine actors and seven fully specified use cases
(Chapter~\ref{chapter:Survey}) realised over a relational schema of thirty-two
tables; the backend contains 176 Java classes with 11,821 lines of source code
and the frontend 10,575 lines of JavaScript/JSX, both measured with the
\texttt{cloc} tool excluding blank and comment lines, and the business rules
with the highest consistency risk are pinned by fifteen automated tests. The
system implements the core modules required by the project scope: master data,
purchasing, goods receipt, inventory, lot management, sales, goods issue,
invoice, accounting, customer payment, delivery, dashboard, authentication, and
user management. Measured on a realistic dataset of approximately eight thousand
transactional rows, the principal read endpoints respond in under sixty
milliseconds on average, and a targeted optimisation of the delivery-listing
endpoint reduced its response time from about 2{,}380~ms to 47~ms
(Section~\ref{sec:perf}).

The key contributions are the FEFO lot-allocation mechanism, the inventory
reservation mechanism at sales order approval, the failed-delivery recovery
procedure, and the role-based workflow enforcement. These contributions address
the most important risks of a distributor: expired stock, overselling,
inconsistent recovery after delivery failure, and uncontrolled access to
approval and warehouse operations.

Several limitations remain, and stating them precisely is more useful than
overstating the result. The system is a prototype validated in a
single-organisation, single-instance setting. Its performance was measured only
under sequential requests against one backend instance
(Section~\ref{sec:perf}); behaviour under high concurrent load has not been
load-tested, and the system has not been evaluated through long-term operation
in a real production environment. On the functional side, although the
accounting module records simplified double-entry journal entries (each with a
debit and a credit account) automatically from source documents and tracks
customer payments, it does not yet produce full financial statements such as
balance sheets and income statements, perform period-end closing, or integrate
with external accounting software, and the system does not yet support automatic
route optimisation, barcode scanning, or demand forecasting.

\section{Future Work}
% GUIDE: first, work needed to complete existing features; then new directions
% (e.g. route optimisation for delivery, demand forecasting, barcode/QR scanning,
% mobile shipper app, multi-warehouse transfer, reporting/BI).
Future work should first complete and harden the existing functions. Although the
rules with the highest consistency risk are already covered by automated tests
(Chapter~\ref{chapter:Design}), the test set should be extended to the flows that
are currently exercised only through manual end-to-end scenarios, namely
goods-receipt confirmation, sales-order cancellation, and failed-delivery
recovery. The failed-delivery mechanism should add an explicit
database-level recovery marker so that repeated calls cannot restore the same
goods issue twice. The performance evaluation should also be extended from
sequential measurement to a concurrent load test that establishes the
throughput and tail latency the system sustains as the number of simultaneous
users grows. The reporting module should be validated with larger sample data.

At the architectural level, the goods-issue path currently serialises concurrent
deductions with pessimistic write locks taken on the affected \texttt{Inventory}
rows, acquired in ascending product-identifier order to prevent deadlock
(Chapter~\ref{chapter:Design}). This guarantees that lot quantities cannot be
lost-updated, but it also turns those rows into a serialisation point that could
become a throughput bottleneck if transaction volume rises sharply. A natural
evolution is to move toward an event-driven architecture in which stock-issue
requests are placed on a message broker such as Apache Kafka or RabbitMQ and
consumed asynchronously \cite{hohpe2003eip, kleppmann2017ddia}. Processing the
issue commands through a queue would relieve the direct lock contention on the
relational database and decouple request acceptance from stock settlement, at
the cost of an eventual-consistency model that the user interface and the
reservation logic would have to accommodate.

The next development direction is to improve warehouse and delivery efficiency.
Barcode or QR scanning can reduce manual input errors during goods receipt,
goods issue, and inventory counting. A mobile-friendly shipper interface or
dedicated mobile application can make delivery updates easier outside the
office. Route optimization can help delivery administrators reduce travel time
and fuel cost when assigning orders to trips. Multi-warehouse transfer should be
added for distributors that move stock between warehouses before shipment.

Another direction is decision support. Demand forecasting can use historical
sales data to suggest reorder quantities. Business intelligence dashboards can
provide deeper analysis of revenue, stock turnover, expired stock, customer
debt, and supplier performance. Finally, the accounting module can be extended
with financial statements, period-end closing, payment reconciliation, and
integration with external accounting software.

A further refinement concerns the integrity of accounting records during
failed-delivery recovery. At present the system cancels the sales invoice
attached to a failed delivery only while that invoice is still in a cancellable
state and no payment has been recorded against it; once a payment exists,
cancellation is withheld and the case is escalated for manual settlement
(Chapter~\ref{chapter:Design}). A more rigorous approach, consistent with
accounting practice, is to treat an issued document as immutable and never
delete it: rather than cancelling the invoice, the system would automatically
generate a credit note that offsets the customer's receivable. This would
preserve a complete audit trail for every issued document, unify the paid and
unpaid cases under a single rule, and align the recovery procedure with the
controls expected of an accounting information system \cite{romney2021ais}.
